\exercisegroup

\begin{enumerate}
\def\labelenumi{\arabic{enumi})}
\tightlist
\item
  设 B 是无穷维希尔伯特空间 E 中的一个标准正交基。
\end{enumerate}

\begin{enumerate}
\def\labelenumi{\alph{enumi})}
\tightlist
\item
  证明：E 中每个处处稠密子集的基数至少等于 B 的基数，且 E 中存在一个与 B 等势的处处稠密集。
\item
  证明 \(\text{Card}(E) = \text{Card}(B^{N})\)（为看出 \(\text{Card}(E) \leqslant \text{Card}(B^{N})\)，利用 a)）。
\item
  证明：若 \(\text{Card}(B) \leqslant \text{Card}(R)\)，则 E 的每个代数基的基数都等于 \(\text{Card}(R) = 2^{\aleph_{0}}\)（利用 II, p.~80, 习题 24, c)）；然而，若 \(\text{Card}(B) > \text{Card}(R)\)，则 E 的每个代数基都与 \(B^{N}\) 等势（利用 b) 及 A, II, § 7, 习题 3, d)）。
\end{enumerate}

¶ 2) a) 设 \(\mathrm{E}_1, \mathrm{E}_2\) 是两个希尔伯特空间，其希尔伯特维数分别是两个基数 \(m\) 与 \(n\)，满足 \(m < n \leqslant m^{\aleph_{0}}\)。设 \(\mathrm{E} = \mathrm{E}_1 \oplus \mathrm{E}_2\) 是 \(\mathrm{E}_1, \mathrm{E}_2\) 的希尔伯特和，并设 \((b_\lambda)_{\lambda \in L}\) 是 \(\mathrm{E}_2\) 的一个标准正交基。证明：存在一个代数无关系 \((a_\lambda)_{\lambda \in L}\)（在 \(\mathrm{E}_1\) 中；cf.~习题 1, c)）；设 \(H\) 是 \(E\) 的由族 \((a_\lambda + b_\lambda)_{\lambda \in L}\)（代数地）生成的子空间。证明 \(\overline{H}\) 的希尔伯特维数等于 \(H\)（注意从 \(H\) 到 \(\mathrm{E}_2\) 上的正交投影是处处稠密的，并利用习题 1, a)）。若 \(S\) 是 \(H\) 的一个标准正交子集，证明 \(S \cap E_2 = \varnothing\)；并由此推出 \(\operatorname{Card}(S) \leqslant m\)（注意 \(E_1\) 的一个标准正交基的每个元素都与 \(S\) 的每个元素正交，至多除去一个可数无穷子集）。

\begin{enumerate}
\def\labelenumi{\alph{enumi})}
\setcounter{enumi}{1}
\tightlist
\item
  设 \(E_{3}\) 是希尔伯特维数为 \(p \geqslant n\) 的希尔伯特空间，并设 F 是希尔伯特和 \(E \oplus E_{3}\)。设 G 是 F 的子空间 \(H + E_{3}\)。证明 \(\overline{G}\) 的希尔伯特维数是 p。若 T 是 G 的一个标准正交子集，证明 \(T \cap (E_{2} + E_{3}) \subset E_{3}\)；并推出 T 到 \(E_{2}\) 上的正交投影的基数至多是 m（按 a) 的方式论证）。由此得出结论：G 没有标准正交基，方法是注意从 G 到 \(E_{2}\) 上的正交投影在 \(E_{2}\) 中处处稠密。
\end{enumerate}

\begin{enumerate}
\def\labelenumi{\arabic{enumi})}
\setcounter{enumi}{2}
\item
  证明：在每个 Hausdorff 且不完备的准希尔伯特空间 E 中，存在一个闭超平面，它在 E 中的正交子空间缩为 0。由此推出：若 E 满足第一可数公理，则 E 中存在一个非完全的标准正交族，它不含于任何标准正交基。
\item
  设 E 是 Hausdorff 准希尔伯特空间，\((\mathrm{E}_{i})_{i\in\mathbf{I}}\) 是 E 的完备向量子空间组成的一个族，按包含关系良序，使得所有 \(E_{i}\) 的并在 E 中处处稠密。证明：E 中存在一个标准正交基 \((e_{\alpha})_{\alpha\in\mathbf{A}}\)，具有下列性质：对每个 \(i\in I\)，所有 \(e_{\alpha}\)（属于 \(E_{i}\) 者）组成的集是 \(E_{i}\) 的一个标准正交基。（考虑 E 中所有满足下列性质的标准正交子集 S 组成的集：对每个 \(i\in I\)，S 中每个不属于 \(E_{i}\) 的向量都与 \(E_{i}\) 正交；并取这个集的一个极大元。）由此推出 V, p.~24 的推论的一个新证明。
\item
  证明：对希尔伯特维数无穷的希尔伯特空间 E，存在一个从 E 到 E 的异于 E 的闭向量子空间上的同构。
\item
  设 E 是希尔伯特空间，\((e_{i})_{i\in I}\) 是 E 的一个标准正交基。证明：若 \((a_{i})_{i\in I}\) 是 E 中的拓扑无关族，满足 \(\sum_{i\in I}\|e_{i}-a_{i}\|^{2}<+\infty\)，则族 \((a_{i})\) 是完全的。（设 J 是 I 的一个有限子集；证明存在一个从 E 到其自身中的连续线性映射 u，使得 \(u(e_{i})=e_{i}\)（对 \(i\in J\)）、\(u(e_{i})=a_{i}\)（对 \(i\notin J\)），且通过适当选取 J 可使 \(u-1_{E}\) 的范数任意小；然后利用 IV, p.~65, 习题 17。）
\item
  给定 \(n\) 个点 \(x_{i} (1 \leqslant i \leqslant n)\)（在一个 Hausdorff 准希尔伯特空间 \(E\) 中），我们把这 \(n\) 个点的 Gram 行列式理解为行列式
\end{enumerate}

\[
\mathrm{G} (x _ {1}, \dots , x _ {n}) = \det (\langle x _ {i} | x _ {j} \rangle).
\]

\begin{enumerate}
\def\labelenumi{\alph{enumi})}
\tightlist
\item
  证明 \(\mathrm{G}(x_{1},\ldots,x_{n})\geqslant0\)，且为使 \(x_{1},\ldots,x_{n}\) 构成一个无关组，必须且只需 \(\mathrm{G}(x_{1},\ldots,x_{n})\neq0\)（假设 \(\dim(\mathrm{E})\geqslant n\)，考虑一个包含 \(x_{1},\ldots,x_{n}\) 的 n 维子空间的正交基）。
\item
  证明：若 \(x_{1},\ldots,x_{n}\) 是 E 中的一个无关组，则一点 \(x\in E\) 到由 \(x_{1},\ldots,x_{n}\) 生成的向量子空间 V 的距离等于 \((\mathrm{G}(x,x_{1},\ldots,x_{n})/\mathrm{G}(x_{1},\ldots,x_{n}))^{1/2}\)（求出 x 在 V 上的正交投影的表达式）。
\item
  设 \((x_{n})\) 是 E 中的点的无穷序列。为使族 \((x_{n})\) 拓扑无关，必须且只需：对每个整数 p\textgreater0，
\end{enumerate}

\[
\sup _ {n} \left(\mathrm{G} \left(x _ {1}, \dots , x _ {p - 1}, x _ {p + 1}, \dots , x _ {n}\right) / \mathrm{G} \left(x _ {1}, \dots , x _ {n}\right)\right) <   + \infty
\]

（利用 b)）。

\begin{enumerate}
\def\labelenumi{\arabic{enumi})}
\setcounter{enumi}{7}
\item
  设 \(\mathbf{E}\) 是具有可数无穷标准正交基 \((e_n)_{n\geq 1}\) 的希尔伯特空间。设 \(\mathbf{A}\) 是 \(\mathbf{E}\) 中由所有点 \(\left(1 - \frac{1}{n}\right)e_n\)（\(n \geqslant 1\)）组成的集的闭凸包。证明：不存在任何点对 \(x, y\)（在 \(A\) 中），其距离等于 \(A\) 的直径（与 IV, p.~54, 习题 12 比较）。
\item
  \begin{enumerate}
  \def\labelenumii{\alph{enumii})}
  \tightlist
  \item
    设 E 是满足第一可数公理的无穷维实希尔伯特空间。设 \((a_{n})_{n\geqslant0}\) 是 E 中的点的一个自由族，使得 \((a_{2n})\) 与 \((a_{2n+1})\) 两族中的每一族都在 E 中是完全的（II, p.~80, 习题 26, a)）。设 F 与 G 是 E 的两个向量子空间，\((a_{2n})\) 与 \((a_{2n+1})\) 分别是它们的（代数）基。空间 F 与 G 由双线性型 \(\langle y|z\rangle\) 赋以分离对偶。证明：若用 B 表示 E 中的单位球，则在赋以拓扑 \(\sigma(\mathrm{F},\dot{\mathrm{G}})\) 的空间 F 中，凸集 \(F\cap B\) 是闭的，但没有任何闭支撑超平面。
  \end{enumerate}
\end{enumerate}

\begin{enumerate}
\def\labelenumi{\alph{enumi})}
\setcounter{enumi}{1}
\tightlist
\item
  设 \((b_n)_{n\geqslant 1}\) 是 \(\mathbf{B}\) 中的一个处处稠密集，并对每个 \(x\in \mathrm{E}\)，令 \(u(x)\) 为序列 \((\langle b_k|x\rangle /k)_{k\geqslant 1}\)。证明 \(u\) 是一个从 \(\mathrm{E}\) 到希尔伯特空间 \(\ell_{\mathbf{R}}^2 (\mathbf{N})\) 中的单射连续线性映射，且 \(u(\mathbf{B})\) 是紧的。证明：在赋范子空间 \(\mathrm{L} = u(\mathrm{F})\)（\(\ell_{\mathbf{R}}^2 (\mathbf{N})\) 的）中，集 \(u(\mathbf{B}\cap \mathbf{F})\) 是闭的、凸的且预紧的，但没有任何闭支撑超平面（注意：若 \(f\) 是 \(\mathbf{L}\) 上的连续线性型，则 \(f\circ u\) 是 \(\mathrm{F}\) 上关于拓扑 \(\sigma (\mathrm{F},\mathrm{G})\) 的连续线性型）。
\end{enumerate}

\begin{enumerate}
\def\labelenumi{\arabic{enumi})}
\setcounter{enumi}{9}
\tightlist
\item
  设 \(\mathbf{E}\) 是满足第一可数公理的无穷维实希尔伯特空间，\((e_n)_{n\geqslant 1}\) 是 \(\mathbf{E}\) 的一个标准正交基。
\end{enumerate}

\begin{enumerate}
\def\labelenumi{\alph{enumi})}
\item
  设 \(\mathbf{A}\) 是由所有点 \(e_n / n\)（在 \(\mathbf{E}\) 中）组成的集的闭均衡凸包。证明 \(\mathbf{A}\) 是紧的，且不存在 \(\mathbf{A}\) 在点 0 处的任何闭支撑超平面，但存在过 0 的直线 \(\mathbf{D}\) 使得 \(\mathbf{D} \cap \mathbf{A} = \{0\}\)。
\item
  设 \(\mathbf{F}\) 是希尔伯特和 \(\mathbf{E} \oplus \mathbf{R}\)，\(e_0\) 是一个与诸 \(e_n\)（\(n \geqslant 1\)）一起构成 \(\mathbf{F}\) 的标准正交基的向量。若 \(\mathbf{B}\) 是 \(\{e_0\} \cup \mathbf{A}\) 的闭凸包，证明：存在一个闭线段 \(\mathbf{L}\)（在 \(\mathbf{F}\) 中，以 0 为中点），使得 \(\mathbf{L} \cap \mathbf{B} = \{0\}\)，但不存在任何过 0 且分离 \(\mathbf{L}\) 与 \(\mathbf{B}\) 的闭超平面（尽管存在 \(\mathbf{B}\) 在 0 处的闭支撑超平面）。
\end{enumerate}

\begin{enumerate}
\def\labelenumi{\arabic{enumi})}
\setcounter{enumi}{10}
\tightlist
\item
  设 \(\mathbf{E}_1, \mathbf{E}_2\) 是两个满足第一可数公理的无穷维实希尔伯特空间，E 是希尔伯特和 \(E_{1} \oplus E_{2}\)（我们将把它与乘积 \(E_{1} \times E_{2}\) 等同）。设 \((e_{n})_{n \geqslant 1}\) 是 \(E_{1}\) 的一个标准正交基；在 \(E_{2}\) 中，设 A 是包含 0 的紧凸集，D 是过 0 的直线，使得 \(D \cap A = \{0\}\)，并假设不存在 A 在点 0 处的闭支撑超平面（习题 10）。设 \((\alpha_{n})\)、\((\beta_{n})\) 是两个由数 \(\geqslant 0\) 组成的序列，满足 \(\lim_{n \to \infty} \beta_{n} = 0\) 且 \(\sum_{n} \alpha_{n}^{-1} < 1\)。设 P 是所有点
\end{enumerate}

\(\sum_{n}\xi_{n}e_{n}\)（属于 \(\mathrm{E}_1\)）组成的集，这些点满足 \(0\leqslant \xi_n\leqslant \alpha_n\)（对每个 \(n\geqslant 1\)）。最后，设 Q 为

E 中由所有点 \((\alpha_{n}e_{n}, x + \beta_{n}a)\) 组成的集的闭凸包，其中 \(n \geqslant 1\)，\(a \neq 0\) 是 D 中一个固定的点，而 \(x\) 取遍 A。

\begin{enumerate}
\def\labelenumi{\alph{enumi})}
\item
  证明 \(\mathbf{P} \cap \mathbf{Q} = \emptyset\)，且 E 中不存在分离 \(\mathbf{P}\) 与 \(\mathbf{Q}\) 的闭超平面。
\item
  设 F 是希尔伯特和 \(E \oplus R\)，并设 c 是 F 中不含于 E 的任意一点。证明分别由 P 与 Q 生成的、以 c 为顶点的尖凸锥 \(P_{1}\)、\(Q_{1}\) 在 F 中是闭的，且 F 中不存在分离 \(P_{1}\) 与 \(Q_{1}\) 的闭超平面（为看出 \(P_{1}\) 与 \(Q_{1}\) 是闭的，证明 P 与 Q 都不含任何半直线）。
\end{enumerate}

\begin{enumerate}
\def\labelenumi{\arabic{enumi})}
\setcounter{enumi}{11}
\item
  设 E 是无穷维实希尔伯特空间。证明：在 E 上存在无穷多个复希尔伯特空间结构，使得 E 是这些复希尔伯特空间底下的实局部凸空间（II, p.~61）。（为证明满足 \(u^{2}(x) = -x\) 的、E 的拓扑向量空间结构的自同构 u 的存在性，利用 E 的一个标准正交基；然后应用 V, p.~60, 习题 1。）给出一个例子，以说明该命题不能推广到 Hausdorff 不完备的准希尔伯特空间（考虑这样一个空间中的一个处处稠密超平面）。
\item
  设 E 是满足第一可数公理的无穷维希尔伯特空间，\((e_{n})_{n\in\mathbf{Z}}\) 是 E 的一个标准正交基，其指标集是有理整数集。用 u 表示 E 到其自身上的等距映射，它满足 \(u(e_{n}) = e_{n+1}\)（对所有 \(n \in Z\)），并置
\end{enumerate}

\[
f (x) = \frac {1}{2} (1 - \| x \|) e _ {0} + u (x).
\]

\begin{enumerate}
\def\labelenumi{\alph{enumi})}
\item
  设 B 是 E 中的单位球，S 是 E 中的单位球面。证明 \(f\) 在 B 上的限制是从 B 到其自身上的一个同胚（注意 \(u\) 在 S 上的限制是从 S 到其自身上的同胚），且不存在任何点 \(x_0 \in \mathbf{B}\) 使得 \(f(x_0) = x_0\)（用 \(x_0\) 关于 \((e_n)\) 的坐标来表示）。
\item
  对每个 \(x \in \mathbf{B}\)，令 \(g(x)\) 为 \(S\) 与以 \(f(x)\) 为原点且过 \(x\) 的半直线的交点。证明 \(g\) 是一个从 \(B\) 到 \(S\) 上的连续映射，使得 \(g(x) = x\)（对所有 \(x \in S\)；与 GT, VI, § 2, 习题 8 比较）。由此推出：存在 \(x_0 \in S\) 以及一个连续映射 \(h\)（从 \(S \times [0,1]\) 到 \(S\) 上），使得 \(h(x,0) = x_0\) 且 \(h(x,1) = x\)（对所有 \(x \in S\)）。
\end{enumerate}

\begin{enumerate}
\def\labelenumi{\arabic{enumi})}
\setcounter{enumi}{13}
\tightlist
\item
  \begin{enumerate}
  \def\labelenumii{\alph{enumii})}
  \tightlist
  \item
    设 \((\mathrm{E}_n)_{n\geqslant 0}\) 是实 Banach 空间的无穷序列，E 是乘积 \(\mathrm{F} = \prod_{n = 0}^{\infty}\mathrm{E}_n\) 的向量子空间，由所有序列 \(x = (x_{n})\)（满足 \(\sum_{n}\| x_{n}\|^{2} < + \infty\)）组成。证明 E 上的函数 \(\| x\| = (\sum_{n}\| x_{n}\|^{2})^{1 / 2}\) 是一个范数，且 E 对这个
  \end{enumerate}
\end{enumerate}

范数是完备的：我们说 E 是 Banach 空间 \(E_{n}\) 的希尔伯特和。

\begin{enumerate}
\def\labelenumi{\alph{enumi})}
\setcounter{enumi}{1}
\item
  证明强对偶 \(\mathbf{E}'\)（\(\mathbf{E}\) 的）可与强对偶 \(\mathrm{E}_n'\)（空间 \(\mathrm{E}_n\) 的）的希尔伯特和等同，且若 \(x' = (x_n') \in \mathrm{E}'\)，则 \(\langle x, x' \rangle = \sum_{n} \langle x_n, x_n' \rangle\)（若 \(u\) 是 E 上的连续线性型，\(u_{n}\) 是它在 \(E_{n}\)（视为 E 的子空间）上的限制，\(a_{n}\) 是 \(E_{n}\) 中满足 \(\|a_{n}\|=1\) 的一点，证明：对每个实数序列 \((\lambda_{n})\)（满足 \(\sum_{n}\lambda_{n}^{2}<+\infty\)），以 \(\lambda_{n}u_{n}(a_{n})\) 为通项的级数收敛，并由此推出 \(\sum_{n}(u_{n}(a_{n}))^{2}<+\infty\)，例如利用关于 \(\ell^{2}(\mathbf{N})\) 的 Banach-Steinhaus 定理）。
\item
  由 \(b)\) 推出：当每个 \(E_n\) 都自反时，\(E\) 是自反的。特别地，若取 \(E_n\) 为空间 \(\mathbf{R}^n\)（赋以范数 \(\| x \| = \sup_{1 \leqslant i \leqslant n} |\xi_i|\)，对 \(x = (\xi_i)_{1 \leqslant i \leqslant n}\)），证明 E 是自反的，但不存在任何与 E 的拓扑相容且使 E 一致凸的范数（V, p.~67, 习题 31）。
\end{enumerate}

¶ 15) a) 对每个整数 n \textgreater{} 0，设 \(a^{(n)}\) 是 IV, p.~63, 习题 8 中定义的二重序列。设 E 是由实数的二重序列 \(x = (x_{ij})\) 组成的向量空间，使得对每个整数 n \textgreater{} 0，有 \(p_n(x) = (\sum_{i,j} a_{ij}^{(n)} |x_{ij}|^2)^{1/2} < +\infty\)。证明诸 \(p_n\) 是 E 上的半范数，

且赋以由这些半范数定义的拓扑的 E 是 Fréchet 空间且是 Montel 空间（按 IV, p.~60, 习题 11 的方式论证）。

\begin{enumerate}
\def\labelenumi{\alph{enumi})}
\setcounter{enumi}{1}
\item
  证明 E 的对偶可与空间 \(E'\) 等同，后者由所有满足下列条件的二重序列 \(x' = (x'_{ij})\) 组成：至少对一个指标 n，有 \(\sum_{i,j} (a_{ij}^{(n)})^{-1} |x'_{ij}|^2 < +\infty\)。
\item
  对每个 \(x = (x_{ij}) \in \mathbf{E}\)，证明 \(\sum_{j=1}^{\infty} \left( \sum_{i=1}^{\infty} |x_{ij}|^2 \right) < +\infty\)（利用 Cauchy-Schwarz 不等式）；对每个 \(j \geqslant 1\)，置 \(y_j = \sum_{i=1}^{\infty} x_{ij}\)；于是序列 \(u(x) = (y_j)\) 属于希尔伯特空间 \(\ell^2(\mathbf{N})\)。证明 \(u\) 是一个从 \(E\) 到 \(\ell^2(\mathbf{N})\) 上的严格满态射；由此推出：\(\ell^2(\mathbf{N})\) 中存在弱紧集，它们不是在 \(u\) 之下 \(E\) 中一个有界集的像（按 IV, p.~63, 习题 8 的方式论证）。
\end{enumerate}

¶ 16) a) 设 \(\Lambda\) 是递增映射 \(\lambda : \mathbf{N} \to \mathbf{R}_{+}^{*}\) 组成的集；对每个整数 \(n \geqslant 0\) 及每个 \(\lambda \in \Lambda\)，令 \(\phi_{n}(\lambda) = \lambda(n)\)。设 E 是所有满足下列条件的映射 \(x: \Lambda \to \mathbf{C}\) 组成的集：对每个 \(n \in \mathbf{N}\)，有 \(p_{n}(x) = \left( \sum_{\lambda \in \Lambda} |x(\lambda)|^{2} \phi_{n}(\lambda) \right)^{1/2} < +\infty\)。证明 E 是一个向量空间，

且诸 \(p_{n}\) 是定义一个自反 Fréchet 空间结构的半范数。

\begin{enumerate}
\def\labelenumi{\alph{enumi})}
\setcounter{enumi}{1}
\item
  设 \(\mathbf{B}\) 是 \(\mathbf{E}\) 中的一个有界集，并令 \(\alpha_{n} = \sup_{x\in \mathbf{B}}p_{n}(x)\)；设 \(\lambda_0\) 是 \(\Lambda\) 的一个元素，使得 \(\lim_{n\to \infty}\lambda_0(n)^{-1}\alpha_n^2 = 0\)。证明 \(x(\lambda_0) = 0\)（对所有 \(x\in \mathbf{B}\)），因而集 \(\mathbf{B}\) 在 \(\mathbf{E}\) 中不是完全的。
\item
  设 \((\mathrm{U}_n)\) 是 E 中的凸均衡零邻域的一个可数基本系；若 \(\mathrm{U}_n^\circ\) 关于 \(\mathrm{E}'\) 上的强拓扑可度量化，则存在 E 中有界集的序列 \((\mathrm{B}_{nm})_{m\geqslant 0}\)，使得诸集 \(\mathbf{B}_{nm}^{\circ}\cap \mathbf{U}_n^\circ\) 构成 \(\mathbf{U}_n^\circ\) 中 0 关于强拓扑的一个基本邻域系。由 \(b\) ) 推出：存在一个整数 \(n\)，使得 \(\mathrm{U}_n^\circ\) 关于强拓扑不可度量化（利用 III, p.~38 的习题 5）。
\end{enumerate}

